% GENERATED FILE. Edit canonical lessons, then run npm run generate:books.
% canonical-source-hash: fnv1a64:bdb75c8ec4c60181
% canonical-lessons: FA-C292-be-nazar-e-man, FA-C292-baraye-hamin, FA-R292-opinion-and-result

\chapter{What I Think, And So}
\label{ch:fa-what-i-think-and-so}
\hlchaptermodality{292}

\begin{chapteropening}
\textbf{By the end of this chapter:} \emph{I can give an opinion with be nazar-e man or fekr mikonam ke, and draw a result from it with bara-ye hamin or pas.}

Complete FA-R292-opinion-and-result: an opinion, then the result drawn from it.
\end{chapteropening}

\section[\texorpdfstring{\fa{به} \fa{نظر} \fa{من} (be nazar-e man)}{به نظر من (be nazar-e man)}]{\fa{به} \fa{نظر} \fa{من} (be nazar-e man) — to my view, which is to say, in my opinion}

\label{lesson:FA-C292-be-nazar-e-man}



\begin{quote}
Say \textquotedblleft{}to seem\textquotedblright{}, then \textquotedblleft{}to think\textquotedblright{}. (\emph{Be nazar residan}, \emph{fekr kardan}.)

A story says what happened. Giving reasons starts one step earlier, with what you think about it.
\end{quote}

\subsection*{You'll want to know: \fa{به} \fa{نظر} \fa{من}}
\textbf{\fa{به} \fa{نظر} \fa{من}} — \emph{be nazar-e man} — \textquotedblleft{}in my opinion\textquotedblright{}, \textquotedblleft{}if you ask me\textquotedblright{}.

\begin{quote}
\textbf{\fa{به} \fa{نظر} \fa{من} \fa{این} \fa{فیلم} \fa{خیلی} \fa{خوب} \fa{است}.} — \emph{be nazar-e man in film kheyli khub ast} — In my opinion this film is very good.
\end{quote}

Word for word it is \textquotedblleft{}to the view of me\textquotedblright{}: \emph{be}, \textquotedblleft{}to\textquotedblright{}, \emph{nazar}, \textquotedblleft{}view\textquotedblright{}, and the ezafe joining \emph{man} as its owner. Change the owner to hand the opinion back: \textbf{\fa{به} \fa{نظر} \fa{شما؟}}, \emph{be nazar-e shomâ?}, \textquotedblleft{}what do you think?\textquotedblright{}.

It is the \emph{nazar} of \emph{be nazar residan}, \textquotedblleft{}to seem\textquotedblright{}, \textquotedblleft{}to reach the view\textquotedblright{}. The present of that verb opens a careful opinion: \textbf{\fa{به} \fa{نظر} \fa{میرسد} \fa{که}...}, \emph{be nazar miresad ke...}, \textquotedblleft{}it seems that...\textquotedblright{}.

\begin{grammarlens}[title={Grammar Lens: fekr mikonam ke}]
The verb you have held since chapter seven, in the present, with \textbf{\fa{که}}, \emph{ke}, \textquotedblleft{}that\textquotedblright{}:

\begin{quote}
\textbf{\fa{فکر} \fa{میکنم} \fa{که} \fa{قطار} \fa{دیر} \fa{میرسد}.} — \emph{fekr mikonam ke qatâr dir miresad} — I think the train will be late.
\end{quote}

\emph{Ke} opens the thought, and the thought keeps its ordinary verb, \emph{miresad}. In speech \emph{ke} is often dropped: \emph{fekr mikonam qatâr dir miresad}.

\textquotedblleft{}I don't think so\textquotedblright{} is the negative with nothing after it: \textbf{\fa{فکر} \fa{نمیکنم}}, \emph{fekr nemikonam}.
\end{grammarlens}

\begin{cousinweb}
\emph{nazar} is borrowed from Arabic \emph{naẓar}, \textquotedblleft{}a look, a view\textquotedblright{} (EtymDB 2.1); \emph{be} and \emph{man} are Persian. So the phrase sets an Arabic noun inside a Persian frame, the way \emph{fekr kardan} sets Arabic \emph{fekr}, \textquotedblleft{}thought\textquotedblright{}, in front of a Persian verb.
\end{cousinweb}

\subsection*{Your turn}
\begin{itemize}
\raggedright
  \item \emph{Say it:} \emph{be nazar-e man in film kheyli khub ast}
  \item \emph{Say it:} hand it back — \emph{be nazar-e shomâ?}
  \item \emph{Say it:} \emph{fekr mikonam ke qatâr dir miresad}, then without \emph{ke}
  \item \emph{Answer:} \emph{qatâr dir miresad?} with \textquotedblleft{}I don't think so\textquotedblright{} — \emph{fekr nemikonam}
\end{itemize}

\begin{tcolorbox}[breakable,colback=teal!4,colframe=teal!35!black,title={Before you move on}]
What does \emph{be nazar-e man} say, word for word? (\textbf{To the view of me.}) \textquotedblleft{}What do you think?\textquotedblright{} (\emph{\textbf{Be nazar-e shomâ?}}) What does \emph{ke} do after \emph{fekr mikonam}? (It \textbf{opens the thought}, and the verb after it stays ordinary.)
\end{tcolorbox}
\section[\texorpdfstring{\fa{برای} \fa{همین} (barâ-ye hamin)}{برای همین (barâ-ye hamin)}]{\fa{برای} \fa{همین} (barâ-ye hamin) — so: the arrow that points forward}

\label{lesson:FA-C292-baraye-hamin}



\begin{quote}
Say \textquotedblleft{}because\textquotedblright{}, two ways. (\emph{Chon}, \emph{zirâ}.) Then \textquotedblleft{}in my opinion\textquotedblright{}. (\emph{Be nazar-e man.})

An opinion needs its reasons, and a reason needs a result. Persian can join the two from either end.
\end{quote}

\begin{grammarlens}[title={Grammar Lens: two arrows}]
One fact, three ways to join it:

\begin{quote}
\textbf{\fa{خسته} \fa{بودم،} \fa{برای} \fa{همین} \fa{زود} \fa{خوابیدم}.} — \emph{khaste budam, barâ-ye hamin zud khâbidam} — I was tired, so I went to sleep early.
\end{quote}

\begin{quote}
\textbf{\fa{چون} \fa{خسته} \fa{بودم،} \fa{زود} \fa{خوابیدم}.} — \emph{chon khaste budam, zud khâbidam} — Because I was tired, I went to sleep early.
\end{quote}

\begin{quote}
\textbf{\fa{زود} \fa{خوابیدم،} \fa{زیرا} \fa{خسته} \fa{بودم}.} — \emph{zud khâbidam, zirâ khaste budam} — I went to sleep early, because I was tired.
\end{quote}

\textbf{\fa{برای} \fa{همین}}, \emph{barâ-ye hamin}, comes after the cause and points forward to the result. \emph{Barâ-ye} is \textquotedblleft{}for\textquotedblright{}, and \emph{hamin} is \textquotedblleft{}this very\textquotedblright{}: \textquotedblleft{}for this very reason\textquotedblright{}. \emph{Chon} and \emph{zirâ} sit on the cause itself and point back to it. \emph{Chon} can open the sentence; \emph{zirâ} mostly follows the result, and sounds written.
\end{grammarlens}

\begin{grammarlens}[title={Grammar Lens: pas, the short one}]
\textbf{\fa{پس}}, \emph{pas}, is \textquotedblleft{}so\textquotedblright{} when you draw a conclusion, often at the front of what you conclude:

\begin{quote}
\textbf{\fa{پس} \fa{قطار} \fa{دیر} \fa{میرسد}.} — \emph{pas qatâr dir miresad} — So the train will be late.
\end{quote}

\emph{Banâbarin}, \textquotedblleft{}therefore\textquotedblright{}, does the same job in writing. \emph{Pas} is also the \emph{pas} of \emph{pas az}, \textquotedblleft{}after\textquotedblright{}: a result comes after its cause.
\end{grammarlens}

\begin{cousinweb}
\emph{hamin} is \emph{ham}, \textquotedblleft{}same\textquotedblright{}, with \emph{in}, \textquotedblleft{}this\textquotedblright{}. EtymDB 2.1 takes \emph{ham} through Middle and Old Persian back to an Indo-European word for \textquotedblleft{}same\textquotedblright{}, the one English kept as \textbf{same}. \emph{Pas} goes back, in EtymDB, through Middle Persian \emph{ps}, read \emph{pas}, to an Indo-European word meaning \textquotedblleft{}behind, afterwards\textquotedblright{}.
\end{cousinweb}

\subsection*{Your turn}
Say these aloud:

\begin{itemize}
\raggedright
  \item \emph{khaste budam, barâ-ye hamin zud khâbidam}
  \item the same, with the cause first and \emph{chon} in front
  \item a conclusion — \emph{pas qatâr dir miresad}
  \item an opinion and its result — \emph{be nazar-e man film khub ast, barâ-ye hamin miravam}
\end{itemize}

\begin{tcolorbox}[breakable,colback=teal!4,colframe=teal!35!black,title={Before you move on}]
Which joiner points forward to the result? (\emph{\textbf{Barâ-ye hamin}}, or \emph{pas}.) Which point back to the cause? (\emph{\textbf{Chon}} and \emph{\textbf{zirâ}}.) What does \emph{hamin} say, taken apart? (\textbf{This very}: \emph{ham} and \emph{in}.)
\end{tcolorbox}
\section[(review)]{Review — what I think, and what follows}

\label{lesson:FA-R292-opinion-and-result}



\begin{quote}
Say \textquotedblleft{}in my opinion\textquotedblright{}, then \textquotedblleft{}I think that\textquotedblright{}. (\emph{Be nazar-e man}, \emph{fekr mikonam ke}.)
\end{quote}

\begin{grammarlens}[title={What you've built}]
Two lessons, and an opinion now has both ends. It opens with \emph{be nazar-e man} or \emph{fekr mikonam ke}, or carefully with \emph{be nazar miresad ke}. Its cause comes with \emph{chon}. Its result comes with \emph{barâ-ye hamin}, or with \emph{pas} when you are drawing a conclusion.

Think back to the winter trip in chapter two hundred and ninety-one, and give your view of it:

\begin{quote}
\textbf{\fa{به} \fa{نظر} \fa{من} \fa{سفر} \fa{خوب} \fa{بود،} \fa{چون} \fa{هتل} \fa{گرم} \fa{بود}.} — \emph{be nazar-e man safar khub bud, chon hotel garm bud} — In my opinion the trip was good, because the hotel was warm.
\end{quote}
\end{grammarlens}

\subsection*{Your turn}
\begin{itemize}
\raggedright
  \item \emph{Say it:} \emph{be nazar-e man safar khub bud, chon hotel garm bud}
  \item \emph{Say it:} the same view with the result after it — \emph{... barâ-ye hamin dobâre miravam}
  \item \emph{Answer:} \emph{be nazar-e shomâ?} — with \emph{fekr mikonam ke} and a reason
  \item \emph{Say it:} \emph{vaqti bâ qatâr bargashtim, khaste budim, barâ-ye hamin zud khâbidim}
\end{itemize}

\begin{tcolorbox}[breakable,colback=teal!4,colframe=teal!35!black,title={Before you move on}]
Which joiner goes on the cause, and which on the result? (\emph{\textbf{Chon}} on the cause; \emph{\textbf{barâ-ye hamin}} on the result.) \textquotedblleft{}I don't think so\textquotedblright{}? (\emph{\textbf{Fekr nemikonam.}}) Which one draws a conclusion in two letters? (\emph{\textbf{Pas.}})
\end{tcolorbox}
